\chapter*{Prefácio} 
\addcontentsline{toc}{chapter}{Prefácio}

Por qual razão o estudante sempre esquece o que estudou? Em particular, no contexto da Física, por que ele se esquece tanto das fórmulas? Por que ele se confunde nelas, não se lembra da sua equação ou de como elas devem ser usadas?

Certa vez aprendi com um professor de psicologia cognitiva --- doutorado pela Universidade de Harvard --- que a memória é o resíduo do pensamento. Ou seja, lembramos daquilo que pensamos -- que questionamos, que discutimos internamente em nossas mentes.

\pullquote[professor de psicologia cognitiva, Harvard]{A memória é o resíduo do pensamento.\\Lembramos daquilo que pensamos.}

Logo, a maneira mais eficiente de se lembrar de alguma coisa é compreendendo o significado daquela coisa: compreender o significado exige que se pense com profundidade sobre o objeto de estudo, que se caminhe por diferentes raciocínios, argumentos e que se conecte pontos aparentemente desconexos para que o significado seja construído.

Dito isso, este é um livro sobre o significado das coisas. Diferente do que aparenta ser o modus operandi padrão do ensino de Física hoje -- danças, músicas e frases engraçadas para que se memorize uma fórmula -- este que vos escreve acredita que o caminho correto é aquele diametralmente oposto, no qual há pouquíssimo esforço em auxiliar o estudante a memorizar a fórmula, mas um esforço colossal em construir, junto dele, o significado -- completo, profundo e amplo -- a respeito daquela fórmula. Acredito, como os cientistas apontam, que se lembrar da fórmula será apenas uma consequência -- direta e imediata -- de trilhar essa jornada.

Passaremos por quase todas as fórmulas de Física do Ensino Médio, discutindo com ampla profundidade de onde e como elas surgem, em quais contextos aparecem, quais são os erros mais comuns que os estudantes tendem a cometer e como elas podem e devem ser utilizadas por eles.

Física não é só fórmulas, é claro --- afinal de contas, Física não é Matemática. Mas ela é feita em linguagem matemática. Sendo assim, este é um livro que escolhe fazer este recorte e discutir, com profundidade, a linguagem da Física: as suas equações matemáticas. Contudo, o leitor que acha que este será um texto puramente matemático haverá de se decepcionar. Como ficará nítido ao longo do texto, as fórmulas, em muitos momentos, serão apenas o nosso ponto de partida para discutir toda a Física contida no pano de fundo que as originou.


% Ao iniciar a empresa de construir com profundidade e de modo inequívoco o significado das fórmulas, este que vos escreve confessa que percebeu ter subestimado, inicialmente, a complexidade de tal tarefa. Sendo assim, para que o leitor possa consumir o material e estudar melhor, o autor concluiu que seria importante dividir tal tarefa em algumas etapas, de modo que o Guia de Fórmulas Explicadas tornou-se uma coleção que se propõe a cumprir tal tarefa em três volumes:

Ao me lançar na empreitada de construir com profundidade e clareza inequívoca o verdadeiro significado das fórmulas, confesso aqui, sem cerimônia, que subestimei a criatura. Imaginei, no começo, que seria tarefa simples, e com uma página para cada equação seria possível dar cabo em tal tarefa. Mas as fórmulas, mesmo aquelas aparentemente mais simples, escondem mundos inteiros atrás de seus símbolos, e explorá-los se mostrou tarefa mais delicada do que parecia à primeira vista.

Por isso, para o leitor poder atravessar esse caminho com mais calma, estudando sem atropelos, concluí que seria sensato dividir o trabalho em etapas, de modo que o Guia de Fórmulas Explicadas acabou se tornando uma coleção organizada em três volumes, cada qual dedicado a cumprir, à sua maneira, essa missão:

\begin{enumerate}
    \item Volume I - Mecânica.
    \item Volume II - Óptica, Ondas e Termodinâmica.
    \item Volume III - Eletromagnetismo.
\end{enumerate}


Convido agora o leitor ao início do nosso passeio físico-matemático. Espero que ele possa se deleitar nessa travessia e usufruir dos mesmos prazeres que este que vos escreve teve ao compreender diversas das coisas que aqui compartilho.





\begin{flushright}
Felipe Guisoli\\
Belo Horizonte, junho de 2026
\end{flushright}